\documentclass[10pt,a4paper]{amsart}
\usepackage{amsmath,amssymb,amsthm,graphicx}
\usepackage{verbatim}
\newenvironment{reference}{\comment}{\endcomment}
%\newenvironment{reference}{}{}
\newenvironment{slogan}{\comment}{\endcomment}
\newenvironment{history}{\comment}{\endcomment}
% For commutative diagrams we use Xy-pic
\usepackage[all]{xy}

% We use 2cell for 2-commutative diagrams.
\xyoption{2cell}
% Theorem environments.
%
\theoremstyle{plain}
\newtheorem{theorem}[subsection]{Theorem}
\newtheorem{proposition}[subsection]{Proposition}
\newtheorem{lemma}[subsection]{Lemma}

\theoremstyle{definition}
\newtheorem{definition}[subsection]{Definition}
\newtheorem{example}[subsection]{Example}
\newtheorem{exercise}[subsection]{Exercise}
\newtheorem{situation}[subsection]{Situation}

\theoremstyle{remark}
\newtheorem{remark}[subsection]{Remark}
\newtheorem{remarks}[subsection]{Remarks}

\numberwithin{equation}{subsection}

% Macros
%
\def\lim{\mathop{\mathrm{lim}}\nolimits}
\def\colim{\mathop{\mathrm{colim}}\nolimits}
\def\Spec{\mathop{\mathrm{Spec}}}
\def\Hom{\mathop{\mathrm{Hom}}\nolimits}
\def\Ext{\mathop{\mathrm{Ext}}\nolimits}
\def\SheafHom{\mathop{\mathcal{H}\!\mathit{om}}\nolimits}
\def\SheafExt{\mathop{\mathcal{E}\!\mathit{xt}}\nolimits}
\def\Sch{\mathit{Sch}}
\def\Mor{\mathop{\mathrm{Mor}}\nolimits}
\def\Ob{\mathop{\mathrm{Ob}}\nolimits}
\def\Sh{\mathop{\mathit{Sh}}\nolimits}
\def\NL{\mathop{N\!L}\nolimits}
\def\CH{\mathop{\mathrm{CH}}\nolimits}
\def\proetale{{pro\text{-}\acute{e}tale}}
\def\etale{{\acute{e}tale}}
\def\QCoh{\mathit{QCoh}}
\def\Ker{\mathop{\mathrm{Ker}}}
\def\Im{\mathop{\mathrm{Im}}}
\def\Coker{\mathop{\mathrm{Coker}}}
\def\Coim{\mathop{\mathrm{Coim}}}

% Boxtimes
%
\DeclareMathSymbol{\boxtimes}{\mathbin}{AMSa}{"02}

%
% Macros for moduli stacks/spaces
%
\def\QCohstack{\mathcal{QC}\!\mathit{oh}}
\def\Cohstack{\mathcal{C}\!\mathit{oh}}
\def\Spacesstack{\mathcal{S}\!\mathit{paces}}
\def\Quotfunctor{\mathrm{Quot}}
\def\Hilbfunctor{\mathrm{Hilb}}
\def\Curvesstack{\mathcal{C}\!\mathit{urves}}
\def\Polarizedstack{\mathcal{P}\!\mathit{olarized}}
\def\Complexesstack{\mathcal{C}\!\mathit{omplexes}}
% \Pic is the operator that assigns to X its picard group, usage \Pic(X)
% \Picardstack_{X/B} denotes the Picard stack of X over B
% \Picardfunctor_{X/B} denotes the Picard functor of X over B
\def\Pic{\mathop{\mathrm{Pic}}\nolimits}
\def\Picardstack{\mathcal{P}\!\mathit{ic}}
\def\Picardfunctor{\mathrm{Pic}}
\def\Deformationcategory{\mathcal{D}\!\mathit{ef}}
\usepackage{cleveref}
\crefname{originalreference}{original reference}{original references}
\Crefname{originalreference}{Original reference}{Original references}
\makeatletter
\newlabel{spaces-descent-lemma-locally-finite-presentation-fppf-local-source}{{spaces-descent-lemma-locally-finite-presentation-fppf-local-source}{source}{}{sourcepointer-1}{}}
\newlabel{spaces-descent-lemma-locally-finite-presentation-fppf-local-source@cref}{{[originalreference][0][]spaces-descent-lemma-locally-finite-presentation-fppf-local-source}{[1][0][]source}{}{}{}}
\makeatother
\title{The middle factor of a composite which is flat and locally of finite presentation is flat and locally of finite presentation}
\author{The Stacks Project Authors}
\date{stacks-project:stacks-morphisms.tex:local-snapshot-20261012; GNU FDL 1.2 or later}
\begin{document}
\maketitle
\noindent\textit{Editorial note.} All delivered mathematical text is the author's own, copied from the pinned Stacks Project chapter source sources/stacks-algebraic-stacks-morphisms/pinned-original-source.tex; no statement or proof step was rewritten, reordered or removed.\par\medskip
\noindent\textit{Editorial note.} Preamble provenance: the Stacks Project keeps one single shared preamble for every chapter; this card takes it directly from the staged sources/stacks-algebraic-stacks-morphisms/preamble.tex as three verbatim spans (lines 8-12, 14-18 and 154-217), so no derived file was written and nothing was added or redefined. The wrapper class and the base packages are supplied by the assembler.\par\medskip
\noindent\textit{Editorial note.} Bibliography provenance: neither the counted statement nor its proof contains a citation (this lane rescans the whole span before writing the card), so the card carries no bibliography and no entry was reconstructed.\par\medskip
\noindent\textit{Editorial note.} Selection provenance: this card was opened under the file-based control contract (CONTROL.md), with the mechanical pre-card gate run on its statement and proof interval before the card was written.\par\medskip
\noindent\textit{Allowed standard prerequisites (editorial).}\begin{itemize}
\item The definition of the properties flat and locally of finite presentation for morphisms of algebraic stacks through smooth covers by algebraic spaces. The author's section on morphisms of finite presentation, in which the counted statement occurs. Reduces the assertion to the corresponding assertion for morphisms of algebraic spaces.
\item The descent of the property of being locally of finite presentation along a surjective flat morphism of finite presentation of algebraic spaces. Invoked at the end of the counted proof (source lines 5870-5873) as a cross-chapter pointer. Supplies the statement for algebraic spaces to which the proof reduces.
\end{itemize}\medskip
\ifdefined\hypertarget\hypertarget{sourcepointer-1}{}\fi
\noindent\textit{Source reference.} Original reference spaces-descent-lemma-locally-finite-presentation-fppf-local-source: sources/stacks-algebraic-stacks-morphisms/pinned-original-source.tex, source lines 3896--3896. Navigation pointer only: the counted text refers to the Stacks item spaces-descent-lemma-locally-finite-presentation-fppf-local-source, which belongs to another chapter of the project and is not part of this card. The pointer carries the target's own label as its number because this chapter's staged snapshot does not contain the other chapter's tag table; its content is not reproduced here.\par\medskip
\begin{lemma}
\label{lemma-flat-finite-presentation-permanence}
Let $\mathcal{X} \to \mathcal{Y} \to \mathcal{Z}$ be morphisms of
algebraic stacks. If $\mathcal{X} \to \mathcal{Z}$ is locally of finite
presentation and $\mathcal{X} \to \mathcal{Y}$ is surjective, flat, and
locally of finite presentation, then $\mathcal{Y} \to \mathcal{Z}$
is locally of finite presentation.
\end{lemma}

\begin{proof}
Choose an algebraic space $W$ and a surjective smooth morphism
$W \to \mathcal{Z}$. Choose an algebraic space $V$ and a surjective smooth
morphism $V \to W \times_\mathcal{Z} \mathcal{Y}$. Choose an algebraic space
$U$ and a surjective smooth morphism $U \to V \times_\mathcal{Y} \mathcal{X}$.
We know that $U \to V$ is flat and locally of finite presentation
and that $U \to W$ is locally of finite presentation.
Also, as $\mathcal{X} \to \mathcal{Y}$ is surjective we see that
$U \to V$ is surjective (as a composition of surjective morphisms).
Hence the lemma reduces to the case of morphisms of algebraic spaces.
The case of morphisms of algebraic spaces is
Descent on Spaces, Lemma
\ref{spaces-descent-lemma-locally-finite-presentation-fppf-local-source}.
\end{proof}

\end{document}
